\cvsection{Expériences professionnelles}

\begin{cventries}
  \cventry
    {Stagiaire ingénieur de recherche en IA}
    {INRIA}
    {Bordeaux, France}
    {Mars 2026 -- Septembre 2026}
    {
      \begin{cvitems}
        \item {Conception et développement en autonomie d’une application fonctionnelle intégrant un agent IA multi-étapes pour la recherche d’articles scientifiques.}
        \item {Implémentation de l’orchestration sans framework dédié avec l’API OpenAI, conception des prompts, gestion des appels d’outils, de l’état et de la mémoire conversationnelle.}
        \item {Développement d’API REST Python/FastAPI documentées avec OpenAPI, stockage dans MongoDB et interface React/Next.js permettant de sélectionner les messages conservés dans le contexte de l’agent.}
        \item {Automatisation de la collecte de métadonnées d’articles et mise en œuvre d’un pipeline RAG avec chunking, embeddings, indexation dans ChromaDB et recherche sémantique.}
        \item {Définition d’un protocole d’évaluation de LLM et comparaison de modèles et de quantifications selon la réussite de la recherche, la latence et la taille du contexte.}
        \item {Déploiement de modèles avec vLLM sur cluster HPC GPU NVIDIA, conteneurisation de composants avec Docker et mise en place d’une CI/CD sur GitHub.}
        \item {Évaluation d’OpenClaw dans un environnement Docker isolé pour analyser les risques liés à l’accès d’un agent aux outils et aux fichiers.}
      \end{cvitems}
    }

  \cventry
    {Stagiaire ingénieur de recherche en apprentissage par renforcement}
    {Université Karunya}
    {Coimbatore, Inde}
    {Juin 2025 -- Septembre 2025}
    {
      \begin{cvitems}
        \item {Développement d’un modèle d’apprentissage par renforcement et de son environnement d’entraînement pour la navigation d’un robot hospitalier, puis déploiement et tests sur Raspberry Pi.}
      \end{cvitems}
    }
\end{cventries}
